\section{Threats to Validity}
\label{sec:threats}

\paragraph{Internal validity.}
\emph{Q2 benchmark artefact.}
For Q2, the ground-truth function (\texttt{sqlite3Select}) is identical to the entry function
specified in the problem prompt; answering the entry function is therefore correct by construction.
On \qwen, the model answered \texttt{sqlite3Select} across 5 of 6 questions in \textsc{NOCTX},
producing a 3 of 3 exact score on Q2.
To prevent this benchmark artefact from skewing conclusions, we exclude Q2 from the primary
comparisons across both models ($n{=}5$) and report the full benchmark ($n{=}6$) separately.
\emph{Forced final answers and turn horizon.}
All 48 baseline tool-condition runs (36 from \qwen{}, comprising 12 full-tool and 24 ablation runs,
and 12 from \nemotron{}) hit the 15-turn cap; every baseline tool answer is
therefore a forced final answer rather than a model-initiated termination.
Our sensitivity probes (Section~\ref{sec:turn_horizon}) confirm that turn truncation was a
limiting factor: extending the budget enabled both tool conditions to solve additional targets
(including Q4, previously unsolved) and prompted earlier autonomous terminations.
\emph{Tool delivery and source granularity confounds.}
The comparison between TOOL and TEXT conditions carries a delivery confound:
tool-condition agents can read raw function bodies via \texttt{get\_function\_source},
whereas text-condition models see only call relationships and signatures in the slice.
Furthermore, \texttt{get\_function\_source} paginates at approximately 100 lines per call;
for large functions such as \texttt{sqlite3Select} ($>$1{,}500 lines), reading the complete body
consumes multiple turns, potentially biasing the agent toward superficial inspection of monolithic functions.
A $(+\text{src})$ text variant injecting full function source was not evaluated.
\emph{Unevaluated graph features.}
None of the six benchmark questions specifically concerns a taint vulnerability or struct field hierarchy bug.
Consequently, the \texttt{get\_type\_hierarchy} and \texttt{get\_taint\_paths} tools, along with the
84 compile-time \texttt{taint\_reaches} edges, were never queried by the benchmark questions.
\emph{Nemotron ablation exclusion.}
When evaluating \nemotron{} on tool-ablation subsets, the NVIDIA API repeatedly returned
empty responses, short non-JSON strings, or malformed outputs across queries.
Because stable ablation runs could not be completed under this API setup, ablation results
are reported for \qwen{} only.
\emph{Quantization and output truncation.}
\qwen was served locally via LM Studio, whereas \nemotron was accessed via hosted inference.
Under \nemotron \textsc{SYN-TEXT} on Q1, the response exceeded the 3{,}000 completion-token limit
and truncated mid-stream, producing an invalid JSON parse (accounting for the 5 of 6 valid JSON count
in Table~\ref{tab:cost}).
\emph{Ground truth confirmation.}
Q3 and Q4 ground truths were confirmed with a project partner on November 11;
there is no independent external validation beyond the course benchmark.

\paragraph{External validity.}
Our study covers one codebase (\sqlite~3.30.1), six benchmark questions, and
two models.
Results may not generalise to other programming languages, codebase scales, bug classes,
or LLM architectures.
\sqlite's function-pointer-heavy architecture means the callback-gap failure
rate may be higher than in object-oriented or functional codebases.

\paragraph{Construct validity.}
The not-in-graph rate is an upper bound on model fabrication because
function-pointer structural gaps also contribute absent edges.
We use \emph{not supported by the graph} rather than \emph{hallucinated} to
reflect this ambiguity.
The 15-turn cap and forced-final mechanism may suppress answers the model would
eventually reach with more turns.
